%%%%%% !TEX program = xelatex x2 + bibtex + xelatex
\documentclass[oneside,doctor]{zjuthesis}
% 默认为单面模式，如需打印请把oneside换为twoside，博士请将master换为doctor
%==============================================================
%==============================================================

%自己需要增加什么 package 或修改什么设置的话，都放在这里吧。
\usepackage{enumerate}
\usepackage[perpage,symbol]{footmisc}
\setfnsymbol{wiley}
\usepackage{hypbmsec}
\usepackage{dsfont}
\usepackage{listings}
\usepackage{xcolor}
%\usepackage{titletoc}
%\usepackage{lineno}
\usepackage{listings}
\usepackage{lstlinebgrd}
%\usepackage{draftwatermark} 
%\SetWatermarkText{DRAFT} %add watermark text 
%\SetWatermarkScale{1} %specify the size of the text
%\SetWatermarkLightness{0.9}
\usepackage{subfigure}
\usepackage[subfigure]{tocloft}

%\renewcommand\cftchapafterpnum{\vskip0pt}
%\renewcommand\cftsecafterpnum{\vskip0pt}
%\renewcommand{\cfttabfont}{\xiaosan\bfseries\Fangsong}
%\usepackage{overpic}
\usepackage{float} % \begin{lstlisting}[float,floatplacement=b]
% Add your keywords here, and have this in a separate file
% and include it in your preamble
\lstdefinestyle{my-fortran-style}{
    language=[90]Fortran,
    basicstyle=\tiny, %\footnotesize, % 字体大小
    numbers=left, 
    numberstyle= \tiny, %\footnotesize, % 编号大小
    keywordstyle= \color{blue!100},
    commentstyle= \color{gray!100}, 
    backgroundcolor = \color{white!10},    
    frame=tb, % 阴影效果 tb, single, shadowbox, trBL, fttt
    rulesepcolor= \color{ red!20!green!20!blue!20} ,
    escapeinside=``, % 英文分号中可写入中文
    xleftmargin=2em,xrightmargin=1.5em, aboveskip=1em,
    framexleftmargin=2.5em,
    escapeinside={\%*}{*)},
    columns=fixed
}
%\lstset{
%    language=[90]Fortran,
%    basicstyle=\tiny, %\footnotesize, % 字体大小
%    numbers=left, 
%    numberstyle= \tiny, %\footnotesize, % 编号大小
%    keywordstyle= \color{blue!100},
%    commentstyle= \color{green!100}, 
%    frame=shadowbox, % 阴影效果
%    rulesepcolor= \color{ red!20!green!20!blue!20} ,
%    escapeinside=``, % 英文分号中可写入中文
%    xleftmargin=2em,xrightmargin=2em, aboveskip=1em,
%    framexleftmargin=2em
%}
\renewcommand{\lstlistlistingname}{程序汇总}
\renewcommand{\lstlistingname}{范例}
%\titlecontents{section}[4em]{}{\contentslabel{2.5em}}{\hspace*{-2em}}{~\titlerule*[0.6pc]{$.$}~\contentspage} % section 目录缩进
%\titlecontents{subsection}[6.5em]{}{\contentslabel{2.5em}}{\hspace*{-2em}}{~\titlerule*[0.6pc]{$.$}~\contentspage} % subsection 目录缩进
\usepackage[titletoc]{appendix}
\usepackage{tikz}
\usetikzlibrary{fit}
\tikzset{every label/.style={font=\small,inner sep=1.5pt}}
\newcommand{\stencilpt}[4][]{\node[circle,fill=white,draw=blue, scale=2.09, label={below:#4},#1] at (#2) (#3) {}}

%\usepackage[hidelinks]{hyperref} % Haowei added for superlink
\usepackage[colorlinks,
            linkcolor=blue,
            anchorcolor=blue,
            citecolor=blue,
            urlcolor=blue,
            bookmarks=true,
            bookmarksnumbered=true%,
            %linktocpage
            ]{hyperref} % Haowei added for superlink
%\hypersetup{linktoc = all}
% \usepackage[super,square]{natbib} % Haowei added
% \renewcommand\refname{参考文献} % Haowei added

\newcommand{\tabincell}[2]{\begin{tabular}{@{}#1@{}}#2\end{tabular}}

\usepackage{emptypage} % method1: make even page empty when start new chapter

%\usepackage{nextpage} % method2: make list start at odd pages, https://latex.org/forum/viewtopic.php?t=21800

\makeatletter
\renewcommand\@pnumwidth{1.5em} % Parts' page numbers greater than 100 are misaligned in my ToC % https://tex.stackexchange.com/questions/381162/parts-page-numbers-greater-than-100-are-misaligned-in-my-toc
\makeatother

%==============================================================
%==============================================================
%\includeonly{chapters/chapter1, chapters/chapter2, chapters/chapter3} %只编译某一章节
\begin{document}
%==============================================================
%==============================================================
%这部分是论文封面、题名页需要的信息，请根据《研究生学位论文编写规则》自行修改
%\raggedbottom %在twoside模式下使用，页面不会强制在最低端拉伸对齐，页面下部有可能留白
%\flushbottom %在twoside模式默认开启，页面会强制在最低端拉伸对齐，段落间距可能增大
%论文分类号
\renewcommand{\zjuclass}{O534}%http://www.ztflh.com/?c=10092

%论文密级
\renewcommand{\zjusecurity}{非涉密}


\renewcommand{\zjutitlec}{XXXXXXXXXXXX}%中文论文题目
\renewcommand{\zjutitlecb}{XXXXXXXXX}%中文论文题目第二行,若无请留空
\renewcommand{\zjutitlecc}{}%中文论文题目第三行,若无请留空
\renewcommand{\zjutitlee}{XXXXXXXXXXXX}%英文论文题目
\renewcommand{\zjutitleeb}{XXXXXXXXX}%英文论文题目第二
\renewcommand{\zjutitleec}{}%英文论文题目第二行,若无请留空
%作者姓名
\renewcommand{\zjuauthor}{XXX}

%作者学号
\renewcommand{\zjuauthorid}{xxxxxxxx}

%指导教师 
\renewcommand{\zjumentor}{XXX ~ 教授}

%合作导师（如果有的话请取消注释）
%\renewcommand{\zjumentorco}{}

%专业名称
\renewcommand{\zjumajor}{等离子体物理}

%研究方向
\renewcommand{\zjusubject}{磁约束聚变}

%所在学院
\renewcommand{\zjuschool}{物理学院}

%提交日期
\renewcommand{\zjuapprovaldate}{202X年X月}

%答辩日期
\renewcommand{\zjudefencedatec}{202X年X月}
\renewcommand{\zjudefencedatee}{Xxxx 2022}

%论文评阅人（格式：姓名 ~ 职称 ~ 单位）
%\renewcommand{\zjurevieweronec}{陈某某 ~ 教授 ~ XX大学}
%\renewcommand{\zjurevieweronee}{Moumou Chen ~Prof. ~XX University}

%\renewcommand{\zjureviewertwoc}{张某某 ~ 教授 ~ XX大学}
%\renewcommand{\zjureviewertwoe}{Moumou Zhang ~Prof. ~XX University}

%\renewcommand{\zjureviewerthreec}{李某某 ~ 教授 ~ XX大学}
%\renewcommand{\zjureviewerthreee}{Moumou Li ~Prof. ~XX University }

%\renewcommand{\zjureviewerfourc}{吴某某 ~ 教授 ~ XX大学}
%\renewcommand{\zjureviewerfoure}{Moumou Wu ~Prof. ~XX University}

%\renewcommand{\zjureviewerfivec}{杨某某 ~ 教授 ~ XX大学}
%\renewcommand{\zjureviewerfivee}{Moumou Yang ~Prof. ~XX University}
\renewcommand{\zjurevieweronec}{}
\renewcommand{\zjurevieweronee}{}

\renewcommand{\zjureviewertwoc}{}
\renewcommand{\zjureviewertwoe}{}

\renewcommand{\zjureviewerthreec}{}
\renewcommand{\zjureviewerthreee}{}

\renewcommand{\zjureviewerfourc}{}
\renewcommand{\zjureviewerfoure}{}

\renewcommand{\zjureviewerfivec}{}
\renewcommand{\zjureviewerfivee}{}

%答辩委员会（姓名\职称\单位）
\renewcommand{\zjucommitteemainc}{}
\renewcommand{\zjucommitteemaine}{}

\renewcommand{\zjucommitteeonec}{}
\renewcommand{\zjucommitteeonee}{}

\renewcommand{\zjucommitteetwoc}{}
\renewcommand{\zjucommitteetwoe}{}

\renewcommand{\zjucommitteethreec}{}
\renewcommand{\zjucommitteethreee}{}

\renewcommand{\zjucommitteefourc}{}
\renewcommand{\zjucommitteefoure}{}

\renewcommand{\zjucommitteefivec}{}
\renewcommand{\zjucommitteefivee}{}


%==============================================================
% 这部分除了“取舍”外，不需要自己修改，必要信息都已在上面设置。
\xiaosi
%封面
%%% 封面

\thispagestyle{empty}

\begin{tabbing}
\Songti\xiaosi 分类号： \= \underline{\makebox[4cm]{\xiaosi\zjuclass}} \= \hspace{4cm} \= \Songti\xiaosi 单位代码：\underline{\makebox[3cm]{10335}} \\
\Songti\xiaosi 密\hspace{4.22mm}级： \> \underline{\makebox[4cm]{\zjusecurity}} \> \> \Songti\xiaosi 学\hspace{8.44mm}号：\underline{\makebox[3cm]{\zjuauthorid}}
\end{tabbing}

\vspace{5mm}

\begin{center}
  \includegraphics[width=52mm]{images/zjuname.pdf}%浙江大学
   %\includegraphics[width=52mm]{images/zjdx}%浙江大学
\end{center}

\ifthenelse{\equal{\zjudegree}{D}}{\centerline{\Songti\xiaoyi 博士学位论文}}{}
\ifthenelse{\equal{\zjudegree}{M}}{\centerline{\Songti\xiaoyi 硕士学位论文}}{}

\vspace{4mm}

\begin{center}
  \includegraphics[width=21mm]{images/zjulogo.pdf}%校徽
  %\includegraphics[width=21mm]{images/standxb.png}%校徽
\end{center}

\vspace{8mm}

\begin{tabbing}
\Songti\sanhao\bfseries 中文论文题目：\=\hspace{-4mm}\xunderline[112mm]{\linespread{1.1}\bfseries\Fangsong\xiaoer\centerline\zjutitlec}
\ifthenelse{\equal\zjutitlecb{}}{}{\\[2mm]\>\hspace{-4mm}\xunderline[112mm]{\linespread{1.1}\bfseries\Fangsong\xiaoer\centerline\zjutitlecb}}
\ifthenelse{\equal\zjutitlecc{}}{}{\\[2mm]\>\hspace{-4mm}\xunderline[112mm]{\linespread{1.1}\bfseries\Fangsong\xiaoer\centerline\zjutitlecc}}
\\[2mm]\Songti\sanhao\bfseries 英文论文题目：\>\hspace{-4mm}\xunderline[112mm]{\linespread{1.1}\bfseries\sanhao\TNR\centerline\zjutitlee}
\ifthenelse{\equal\zjutitleeb{}}{}{\\[2mm]\>\hspace{-4mm}\xunderline[112mm]{\linespread{1.1}\bfseries\sanhao\TNR\centerline\zjutitleeb}}
\ifthenelse{\equal\zjutitleec{}}{}{\\[2mm]\>\hspace{-4mm}\xunderline[112mm]{\linespread{1.1}\bfseries\sanhao\TNR\centerline\zjutitleec}}
\end{tabbing}

\vspace{15mm}

\begin{tabbing}
\hspace{25mm} \Songti\sihao 申 \hspace{-2mm} \= \Songti\sihao 请人姓名： \= \underline{\makebox[7cm]{\sihao\zjuauthor}} \\[2mm]
              \> \Songti\sihao 指导教师： \> \underline{\makebox[7cm]{\sihao\zjumentor}} \\[2mm]
              %\> \Songti\sihao合作导师： \> \underline{\makebox[6cm]{\sihao\zjumentorco}} \\[2mm]
              \> \Songti\sihao 专业名称： \> \underline{\makebox[7cm]{\sihao\zjumajor}} \\[2mm]
              \> \Songti\sihao 研究方向： \> \underline{\makebox[7cm]{\sihao\zjusubject}} \\[2mm]
              \> \Songti\sihao 所在学院： \> \underline{\makebox[7cm]{\sihao\zjuschool}}
\end{tabbing}

\vspace{10mm}

\begin{tabbing}
\hspace{30mm} \= \Songti\xiaosan\bfseries 论文提交日期： \= \underline{\makebox[49mm]{\Songti\xiaosan\bfseries\zjuapprovaldate}}
\end{tabbing}

\ifthenelse{\equal{\zjuside}{T}}{%
  %\newpage\mbox{}%
  \thispagestyle{empty}}{}

%中文题名页
%%% 中文提名页

\newpage
\cleardoublepage  % make list start at odd pages
\thispagestyle{empty}

\vspace{7.5mm}%\vspace{5mm}

\begin{center}
\xunderline[115mm]{\linespread{1.1}\xiaoer\Fangsong\bfseries\centerline\zjutitlec}
\ifthenelse{\equal\zjutitlecb{}}{}{\\[2mm]\xunderline[115mm]{\linespread{1.1}\xiaoer\Fangsong\bfseries\centerline\zjutitlecb}}
\ifthenelse{\equal\zjutitlecc{}}{}{\\[2mm]\xunderline[115mm]{\linespread{1.1}\xiaoer\Fangsong\bfseries\centerline\zjutitlecc}}
\end{center}

\vspace{7.5mm}%\vspace{5mm}

\begin{center}
  \includegraphics[width=21mm]{images/zjulogo.pdf}%校徽
  %\includegraphics[width=21mm]{images/standxb.png}%校徽
\end{center}

\vspace{4.5mm}%\vspace{3mm}

\begin{tabbing}
\hspace{30mm} \= \Songti\xiaosan\bfseries 论文作者签名： \= \underline{\makebox[5cm]{}} \\[8mm]
              \> \Songti\xiaosan\bfseries 指导教师签名： \> \underline{\makebox[5cm]{}}
\end{tabbing}

\vspace{12mm}%\vspace{8mm}

\begin{tabbing}
\hspace{10mm} \Songti\sihao 论文 \hspace{-2.6mm} \= \Songti\sihao 评阅人~1： \= \underline{\makebox[9cm]{\sihao\zjurevieweronec}} \\
              \> \Songti\sihao 评阅人~2： \> \underline{\makebox[9cm]{\sihao\zjureviewertwoc}} \\
              \> \Songti\sihao 评阅人~3： \> \underline{\makebox[9cm]{\sihao\zjureviewerthreec}} \\
              \> \Songti\sihao 评阅人~4： \> \underline{\makebox[9cm]{\sihao\zjureviewerfourc}} \\
              \> \Songti\sihao 评阅人~5： \> \underline{\makebox[9cm]{\sihao\zjureviewerfivec}}
\end{tabbing}

\vspace{12mm}%\vspace{8mm}

\begin{tabbing}
\hspace{5mm}\Songti\sihao 答辩委员 \hspace{-2.6mm} \= \Songti\sihao 会主席： \= \underline{\makebox[9cm]{\sihao\zjucommitteemainc}} \\
          \>    \Songti\sihao~委员~1： \> \underline{\makebox[9cm]{\sihao\zjucommitteeonec}} \\
          \>    \Songti\sihao~委员~2： \> \underline{\makebox[9cm]{\sihao\zjucommitteetwoc}} \\
          \>    \Songti\sihao~委员~3： \> \underline{\makebox[9cm]{\sihao\zjucommitteethreec}} \\
          \>    \Songti\sihao~委员~4： \> \underline{\makebox[9cm]{\sihao\zjucommitteefourc}} \\
          \>    \Songti\sihao~委员~5： \> \underline{\makebox[9cm]{\sihao\zjucommitteefivec}}
\end{tabbing}

\vspace{12mm}%\vspace{8mm}

\begin{tabbing}
\hspace{34mm} \= \Songti\sihao 答辩日期： \= \underline{\makebox[5cm]{\Songti\sihao\zjudefencedatec}} \\
\end{tabbing}

\ifthenelse{\equal{\zjuside}{T}}{%
  %\newpage\mbox{}%
  \thispagestyle{empty}}{}

%英文题名页
%%% 英文标题页

\newpage
\cleardoublepage  % make list start at odd pages
\thispagestyle{empty}

\vspace{10mm}%\vspace{5mm}

\begin{center}
\xunderline[115mm]{\linespread{1.1}\xiaoer\TNR\bfseries\centerline\zjutitlee}
\ifthenelse{\equal\zjutitleeb{}}{}{\\[2mm]\xunderline[115mm]{\linespread{1.1}\xiaoer\Fangsong\bfseries\centerline\zjutitleeb}}
\ifthenelse{\equal\zjutitleec{}}{}{\\[2mm]\xunderline[115mm]{\linespread{1.1}\xiaoer\Fangsong\bfseries\centerline\zjutitleec}}
\end{center}

\vspace{5.5mm}%\vspace{5mm}

\begin{center}
  \includegraphics[width=21mm]{images/zjulogo.pdf}%校徽
  %\includegraphics[width=21mm]{images/standxb.png}%校徽
\end{center}

\vspace{-3mm}%\vspace{-6mm}

\begin{tabbing}
\hspace{20mm} \= \sanhao\bfseries Supervisor's signature: \= \underline{\makebox[5cm]{}}\kill \\
              \> \hspace{7mm} \sanhao\bfseries Author's signature: \> \underline{\makebox[5cm]{}} \\[9mm]
              \> \sanhao\bfseries Supervisor's signature: \> \underline{\makebox[5cm]{}}
\end{tabbing}

\vspace{8mm}%\vspace{8mm}

\begin{tabbing}
\sihao External Reviewers: \= \underline{\makebox[10.8cm]{\sihao\zjurevieweronee}} \\
                    \> \underline{\makebox[10.8cm]{\sihao\zjureviewertwoe}} \\
                    \> \underline{\makebox[10.8cm]{\sihao\zjureviewerthreee}} \\
                    \> \underline{\makebox[10.8cm]{\sihao\zjureviewerfoure}} \\
                    \> \underline{\makebox[10.8cm]{\sihao\zjureviewerfivee}} \\[12mm]%[8mm]
\sihao Examining Committe Chairperson: \= \\
\hspace{41mm} \= \underline{\makebox[10.8cm]{\sihao\zjucommitteemaine}} \\
\sihao Examining Committe Members: \= \\
\hspace{41mm} \= \underline{\makebox[10.8cm]{\sihao\zjucommitteeonee}} \\
             \> \underline{\makebox[10.8cm]{\sihao\zjucommitteetwoe}} \\
             \> \underline{\makebox[10.8cm]{\sihao\zjucommitteethreee}} \\
             \> \underline{\makebox[10.8cm]{\sihao\zjucommitteefoure}} \\
             \> \underline{\makebox[10.8cm]{\sihao\zjucommitteefivee}}
\end{tabbing}

\vspace{3mm}%\vspace{6mm}

\begin{tabbing}
\hspace{28mm} \= \sihao Date of oral defence: \underline{\makebox[5cm]{\sihao\zjudefencedatee}}
\end{tabbing}

\ifthenelse{\equal{\zjuside}{T}}{%
  %\newpage\mbox{}%
  \thispagestyle{empty}}{}
 % 硕士论文请根据需要取舍。
%独创性声明
\input{chapters/announce}

%\renewcommand{\cftpnumalign}{r} % right aligned page numbers
\setlength{\cftfigindent}{0pt} % left aligned fig entries
\setlength{\cfttabindent}{0pt} % left aligned fig entries
\renewcommand{\cftloftitlefont}{\xiaosan\bfseries\Fangsong}
\renewcommand{\cftlottitlefont}{\xiaosan\bfseries\Fangsong}
\renewcommand{\cfttoctitlefont}{\xiaosan\bfseries\Fangsong}
\renewcommand*{\addvspace}[1]{}%single space within entries and double space between
\setlength{\cftbeforechapskip}{0cm}
\setlength{\cftbeforesecskip}{0cm}
\setlength{\cftbeforesubsecskip}{0cm}
%\renewcommand\cftchapdotsep{\cftdotsep}
%\renewcommand\cftchapleader{\bfseries\cftdotfill{\cftchapdotsep}}
\renewcommand{\cftchapfont}{\normalfont\bfseries\Fangsong}% titles in bold
\renewcommand{\cftchappagefont}{\normalfont\bfseries}% page numbers in bold
\renewcommand{\cftdotsep}{5}
\renewcommand{\cftchapleader}{\bfseries\cftdotfill{\cftsecdotsep}}% dot leaders in bold

\frontmatter
\pagenumbering{Roman}

%勘误页
%\input{errata}  % 请根据需要取舍。
%致谢
%%致谢
%\chapter*{\centerline{致\quad 谢}}
\chapter*{致\quad 谢}
\chaptermark{致\quad 谢}
\addcontentsline{toc}{chapter}{致谢}
%\vspace{1cm}

感谢。

\newpage
%序言
%\input{preface} % 请根据需要取舍。
%中文摘要
%% 中文摘要
%\chapter*{\centerline{摘\quad 要}}
\chapter*{摘\quad 要}
\chaptermark{摘\quad 要}
\addcontentsline{toc}{chapter}{摘要}

\vspace{1em}
XXXXX

\vspace{1em}

\noindent{\bfseries\Fangsong 关键词：}XXX，XXXX，XXX


%英文摘要
%% 英文摘要
%\chapter*{\centerline{Abstract}}
\chapter*{Abstract}
\chaptermark{Abstract}
\addcontentsline{toc}{chapter}{Abstract}

\vspace{1em}
XXXXXXX

\vspace{1em}

\noindent{\textbf{Keywords:}} XXX,XXX,XXX

%% 缩写
%\chapter*{\centerline{缩\quad 写}}
\chapter*{缩\quad 写}
\chaptermark{缩\quad 写}
\addcontentsline{toc}{chapter}{缩写}

\vspace{1em}

\noindent
\textbf{AE}: Alfvén Eigenmode, 阿尔芬本征模；\\

 % suoxie

%==============================================================
%这部分不需要自己修改。

%插图和附表清单

\newpage
\cleardoublepage % make list start at odd pages
\phantomsection
\addcontentsline{toc}{chapter}{表目录}
%\cleartooddpage % \usepackage{nextpage}, make list start at odd pages, https://latex.org/forum/viewtopic.php?t=21800
\listoftables
%\chaptermark{表目录}

%\addcontentsline{toc}{chapter}{图目录}
%\listoffigures
%\chaptermark{图目录}
\newpage
\cleardoublepage  % make list start at odd pages
\phantomsection
\addcontentsline{toc}{chapter}{图目录}
%\cleartooddpage % \usepackage{nextpage}, make list start at odd pages, https://latex.org/forum/viewtopic.php?t=21800
\listoffigures
%\chaptermark{图目录}

%术语表
% \printnomenclature
% \chaptermark{术语表}
%==============================================================

%目次页
\newpage
\cleardoublepage  % make list start at odd pages
\phantomsection
%\cleartooddpage % \usepackage{nextpage}, make list start at odd pages, https://latex.org/forum/viewtopic.php?t=21800
\pdfbookmark{\contentsname}{toc}
\tableofcontents
% \addcontentsline{toc}{chapter}{目录} % Haowei added, this has dot, page num, and superlink.
%\addtocontents{toc}{\noindent{目录}\par} % Haowei added, this has no dot, page num, and superlink.
\addtocontents{toc}{\protect\chaptermark{目\quad 次}}
% \addtocontents{toc}{\protect\contentsline {chapter}{\protect\makebox[\linewidth]{目录\hfill}\vspace{-2em}}{}} % Haowei annoted, this cannot work well with hyperref together
\mainmatter

%==============================================================
%\linenumbers % 开始添加行号
\setlength\arraycolsep{2pt} % 设置eqnarray中‘=’前后的空距
\chapter{绪论}\label{chap1}

\section{XXXXX}
能源问题一直是XXXXXXXXX\cite{cowley2016quest}。

\section{XXXXX}

XXXXXX反应方程式如下：
\begin{eqnarray}
    {\rm ^2_1D+^3_1T\rightarrow ^4_2He\left(3.52MeV\right)+^1_0n\left(14.06MeV\right)}, \label{EQDHreaction}
\end{eqnarray}
XXXXXXX：
\begin{eqnarray}
        &&{\rm ^1_0n+^6_3Li\rightarrow ^3_1T+^4_2He}, \label{EQnLi6Reaction}\\
        &&{\rm ^1_0n+^7_3Li\rightarrow ^3_1T+^4_2He+^1_0n}. \label{EQnLi7Reaction}
\end{eqnarray}
 % introduction
\chapter{第二章}\label{chap2}

\section{第一节} \label{WaveAndParticles}
\subsection{第一小节} \label{EPdrift}
如图\ref{figHeidbrink2008Particles}和表\ref{tableExample}所示
\begin{figure}[htp]
	\centering
	\includegraphics[width=0.45\textwidth]{./images/Heidbrink2008Particles.jpg}
	\caption{DIII-D托卡马克中中性束粒子的不同轨道类型相空间分布示意图\cite{Heidbrink2008POP}。}
	\label{figHeidbrink2008Particles}
\end{figure}


\begin{table}[htb]
    \centering
    \caption{XXXXXXXX。}
    \label{tableExample}
      \begin{tabular}{cccccc}
        \toprule
        \multicolumn{1}{m{20mm}}{\heiti\centering AAAA} & \multicolumn{1}{m{20mm}}{\heiti\centering BBBB} & \multicolumn{1}{m{20mm}}{\heiti\centering CCCC} & \multicolumn{1}{m{20mm}}{\heiti\centering DDDD} & \multicolumn{1}{m{20mm}}{\heiti\centering EEEE} & \multicolumn{1}{m{20mm}}{\heiti\centering FFFF} \\
        \midrule
        XXXX   & XXXX & XXXX & XXXX & XXXX & XXXX \\ 
		XXXX   & XXXX & XXXX & XXXX & XXXX & XXXX \\ 
        \midrule
		XXXX   & XXXX & XXXX & XXXX & XXXX & XXXX \\ 
        \bottomrule
      \end{tabular}
  \end{table} % literature
\include{chapters/chapter3} % code

\chapter{第四章}\label{chap4}
 % EP + MHD
\chapter{第五章}\label{chap5}
 % TM + TAE

\chapter{总结和展望}\label{chap6}
\section{总结}

\section{展望}
 % summary
%\addtocontents{toc}{\noindent{附录}\par} % Haowei added, this has no dot, page num, and superlink. I want it show after chapter6, but I donot know why I have to put it before the chapter6???????
%==============================================================
%这也是个不需要自己修改的部分。
%\appendix %设置附录, 章节格式设置, refer to zjuthesis.cls
\titleformat{\chapter}[hang]{\xiaosan\bfseries\Fangsong}{附录\,\thechapter}{1em}{}{}%小三号仿宋加粗，章定义
%\appendixpageoff
\begin{appendices}
\chapter{XXXXXXXX}\label{chapAppendixA}
\chaptermark{附录A\quad XXXXXXXX}
 % CLT与CLT-K代码的GPU移植介绍
\chapter{XXXXXXXX}\label{chapAppendixB}
\chaptermark{附录B\quad XXXXXXXX}
 % CLT代码的实C验位型拓展和共振磁扰动模块
\chapter{XXXXXXXX}\label{chapAppendixC}
\chaptermark{附录C\quad XXXXXXXX}
 % CLT-K代码的电流和压强耦合模型拓展以及等价性校验
\end{appendices}
\backmatter %结束章节自动编号

%参考文献（习惯使用bibtex的可以修改）
\clearpage
\phantomsection
\addcontentsline{toc}{chapter}{参考文献} % 解决目录中没有相应的参考文献的条目问题
\chaptermark{参考文献}
% \input{chapters/bib}
%\setlength\bibindent{200em}
%\setlength\bibhang{0pt}
%\setlength{\parindent}{-200in}
%\setlength{\leftskip}{200in}
%\noindent
%\setlength{\parindent}{100pt}
\makeatletter %https://tex.stackexchange.com/questions/21308/left-aligning-latex-bibliography-display
\renewcommand{\@biblabel}[1]{[#1]\hfill} %https://tex.stackexchange.com/questions/21308/left-aligning-latex-bibliography-display
\makeatother %https://tex.stackexchange.com/questions/21308/left-aligning-latex-bibliography-display
%\renewcommand\bibpreamble{\vspace{1.1\baselineskip}} % choose a suitable vert. skip
\bibliography{chapters/ref} % Haowei added

%附录
%  \include{chapters/appe}

% 出版物
% \input{chapters/publication}

%致谢
%致谢
%\chapter*{\centerline{致\quad 谢}}
\chapter*{致\quad 谢}
\chaptermark{致\quad 谢}
\addcontentsline{toc}{chapter}{致谢}
%\vspace{1cm}

感谢。


% 作者简历
%\chapter*{\centerline{作者简历}}
\chapter*{作者简历}
\chaptermark{作者简历}
\addcontentsline{toc}{chapter}{作者简历}
%~\\[4pt]
\noindent {\Kaiti\textbf{XXX，男，199X年出生，X族，XXXX人。}}\\[4pt]

\noindent {\Kaiti\textbf{1.教育经历：}}\\[4pt]
\quad\romannumeral1. 201X-201X年，浙江大学，物理学系，物理学专业，理学学士学位；\\[4pt]
\quad\romannumeral2. 201X-202X年，浙江大学，物理学院，等离子体物理专业，攻读理学博士学位。\\[4pt]

\noindent {\Kaiti\textbf{2.攻读博士学位期间发表的论文：}}\\[4pt]
\quad\romannumeral1. {title}, \textbf{my name}, coAuthors' names, Journal name, year, Volume (issue), page;\\[4pt]

\newpage\null\thispagestyle{empty}\newpage
%==============================================================
%==============================================================
\end{document}
